\documentclass[10pt,a4paper]{amsart}
\usepackage[margin=19mm]{geometry}
\usepackage{amsmath,amssymb,mathtools}
\usepackage[scr=rsfs,cal=euler]{mathalfa}
\usepackage{xfrac}
\usepackage[unicode,hidelinks,bookmarks=false]{hyperref}
\newcommand{\ie}{i.e.\ }
\newcommand{\cf}{cf.\ }
\newcommand{\resp}{respectively\ }
\newcommand{\Subsection}{\subsection}
\providecommand{\nobreakdash}{\nobreak\hbox{-}}
\setlength{\emergencystretch}{2em}
\multlinegap0pt
\usepackage{amssymb}
\usepackage{xcolor}
\usepackage{mathtools}
\usepackage[shortlabels]{enumitem}
\usepackage{float}
\usepackage{bm}
\newtheorem{cor}{Corollary}
\title{A note on the multiple zeta functions and their variants at identical arguments}
\author{Pawan Singh Mehta}
\date{Source: arXiv:2603.27266v1 (CC BY 4.0)}
\begin{document}
\maketitle

\noindent\emph{Source and licence.} A note on the multiple zeta functions and their variants at identical arguments. Version arXiv:2603.27266v1 (CC BY 4.0), distributed by arXiv under the Creative Commons Attribution 4.0 International licence, http://creativecommons.org/licenses/by/4.0/, as recorded in the arXiv licence metadata for this exact version and shown on the abstract page https://arxiv.org/abs/2603.27266v1. Full licence notice: \texttt{licenses/arxiv-2603.27266-CC-BY-4.0.txt}.\par\smallskip

\noindent\textit{Editorial scope.} Counted unit: the article's cor cor-value-zero (source0.tex lines 572--579). The article's complete original proof of it is delivered with it (source0.tex lines 581--612, one \texttt{proof} environment of 32 lines). No mathematics is written, repaired or re-derived: the statement and the proof are the article's own lines, in the article's own notation, and no retained line is substituted or re-flowed.  Retained uncounted as the source-local context the proof needs: lines 379--381.  Nothing was omitted from any retained span, so the package carries no omission record.  The retained lines cite 1 works, whose entries are transcribed verbatim from the article's own bibliography source in the e-print and delivered with short editorial labels recorded in the item's reference mapping.  No retained line contains a figure environment, an image-inclusion command or drawing code, so no image is delivered and the package declares no asset dependency.  Editorial formatting only, in addition: an AMS \texttt{amsart} preamble that restates the article's own declarations and macros used by the retained lines, namely the environments \texttt{cor} on separate counters, together with the standard packages those lines need.  The retained environments print in this document as Corollary 1 in order of appearance, whereas the article prints the same items with the numbering of its own full document.  The complete native source of the article as distributed in the arXiv e-print for this exact version is bundled unchanged as \texttt{sources/197334/source0.tex}.
\begin{align}\label{eq-zeta}
\zeta_r(s)=\frac{(-1)^r}{r!}\mathbf{Y}_r\left(-0!\zeta(s), -1!\zeta(2s),\ldots,-(r-1)!\zeta(rs)\right).
\end{align}

\begin{cor}\label{cor-value-zero}
Let $r\geq 1$ be an integer. Then  
\begin{align}\label{eq-zero}
\zeta_r(0)=\frac{(-1)^r}{4^r}\binom{2r}{r}\ \text{ and } \
\zeta^{\star}_r(0)=-\frac{1}{r2^{2r-1}}\binom{2r-2}{r-1}.
\end{align}
%{\rm{\bf{Comment}:} see if better bound can be found as striling numbers are well studied.}
\end{cor}

\begin{proof}
Using \eqref{eq-zeta}, we have
\begin{align*}
\zeta_{r}(0)=\frac{(-1)^r}{r!}{\mathbf{Y}_r\left(-0!\zeta(0), -1!\zeta(0),\ldots,-(r-1)!\zeta(0)\right)}
=\frac{(-1)^r}{r!}\mathbf{Y}_r\left(\frac{1}{2}, \frac{1!}{2}, \ldots, \frac{(r-1)!}{2}\right).
\end{align*}
By definition,
$$
\mathbf{Y}_r\left(\frac{0!}{2}, \frac{1!}{2}, \ldots, \frac{(r-1)!}{2}\right)=\sum_{k=1}^{r}\mathbf{B}_{r,k}\left(\frac{0!}{2}, \frac{1!}{2},\ldots, \frac{(r-1)!}{2}\right)=\sum_{k=1}^{r}\frac{\mathbf{B}_{r,k}\left({0!}, {1!},\ldots, {(r-1)!}\right)}{2^k}.
$$
It is well known that (see \cite[Chapter 5]{L})
$$
\mathbf{B}_{r,k}\left({0!}, {1!},\ldots, {(r-1)!}\right)=|s(r,k)|,
$$
where $|s(r,k)|$ is the unsigned Stirling number of the first kind given by rising factorial, i.e.,
\begin{align}\label{gen-stirling}
x\left(x+1\right)\cdots \left(x+r-1\right)=\sum_{k=1}^{r}{|s(r,k)|}{x^k}.
\end{align}
Therefore, we have
$$
\mathbf{Y}_r\left(\frac{0!}{2}, \frac{1!}{2}, \ldots, \frac{(r-1)!}{2}\right)=\sum_{k=1}^{r}\frac{|s(r,k)|}{2^k}=\prod_{k=0}^{r-1}\left(k+\frac{1}{2}\right),
$$
and
$$
\mathbf{Y}_r\left(-\frac{0!}{2}, -\frac{1!}{2}, \ldots, -\frac{(r-1)!}{2}\right)=\sum_{k=1}^{r}\frac{(-1)^k|s(r,k)|}{2^k}=\prod_{k=0}^{r-1}\left(k-\frac{1}{2}\right),
$$
Hence, we get
$$
\zeta_r(0)=\frac{(-1)^r}{r!}\prod_{k=0}^{r-1}\left(k+\frac{1}{2}\right), \quad \zeta^{\star}_r(0)=\frac{1}{r!}\prod_{k=0}^{r-1}\left(k-\frac{1}{2}\right).
$$
Simplifying the above identities, we get the desired formulae.
\end{proof}

\begin{thebibliography}{99}
\bibitem[L]{L} L. Comtet, {\it Advanced Combinatorics, The Art of Finite and Infinite Expansions}, D. Reildel Publishing Company 1974.
\end{thebibliography}
\end{document}
